\section*{Capítulo \ref{ch:b2:plant-life-cycles} --- Ciclos de vida e reprodução das plantas terrestres}

\begin{solution}{exo:b2:plant-life-cycles:1}
\omterm{def:b2:plant-life-cycles:cycles}{Gametófito}: a geração \omterm{def:b2:unicellular-diversity:unicellular}{multicelular} \omterm{def:b2:meiosis-heredity:meiosis}{haploide}, formada por mitose a partir
de um esporo, que produz gametas por mitose. \omterm{def:b2:plant-life-cycles:cycles}{Esporófito}: a geração
\omterm{def:b2:meiosis-heredity:meiosis}{diploide}, formada por mitose a partir de um zigoto, que produz esporos
por \omterm{def:b2:meiosis-heredity:meiosis}{meiose}. Esporo: uma célula \omterm{def:b2:meiosis-heredity:meiosis}{haploide} produzida por \omterm{def:b2:meiosis-heredity:meiosis}{meiose} que se
desenvolve sem se fundir. Gameta: uma célula \omterm{def:b2:meiosis-heredity:meiosis}{haploide} produzida por
mitose (nas plantas) que precisa se fundir com outra.
\end{solution}

\begin{solution}{exo:b2:plant-life-cycles:2}
As células da folha são $2n = 1260$: esporo 630, célula do \omterm{def:b2:plant-life-cycles:fern}{prótalo}
630, anterozoide 630, oosfera 630, zigoto 1260.
\end{solution}

\begin{solution}{exo:b2:plant-life-cycles:3}
Almofada verde: \omterm{def:b2:plant-life-cycles:cycles}{gametófito}. Haste e cápsula: \omterm{def:b2:plant-life-cycles:cycles}{esporófito}. Fronde de
\omterm{def:b2:plant-life-cycles:fern}{samambaia}: \omterm{def:b2:plant-life-cycles:cycles}{esporófito}. \omterm{def:b2:plant-life-cycles:fern}{Prótalo}: \omterm{def:b2:plant-life-cycles:cycles}{gametófito}. Grão de \omterm{def:b2:plant-life-cycles:heterospory}{pólen}: \omterm{def:b2:plant-life-cycles:cycles}{gametófito}
masculino. Reserva de alimento da semente (num pinheiro): \omterm{def:b2:plant-life-cycles:cycles}{gametófito}
feminino. Embrião da semente: o \omterm{def:b2:plant-life-cycles:cycles}{esporófito} seguinte.
\end{solution}

\begin{solution}{exo:b2:plant-life-cycles:4}
Haplonte (\emph{Chlamydomonas}): \omterm{def:b2:meiosis-heredity:meiosis}{meiose} no zigoto, imediatamente.
Diplonte (animais, \emph{Fucus}): a \omterm{def:b2:meiosis-heredity:meiosis}{meiose} produz os gametas.
Haplodiplonte (\omterm{def:b2:plant-life-cycles:moss}{musgos}, \omterm{def:b2:plant-life-cycles:fern}{samambaias}, plantas com sementes, muitas
algas): a \omterm{def:b2:meiosis-heredity:meiosis}{meiose} produz esporos nos \omterm{def:b2:plant-life-cycles:fern}{esporângios} do \omterm{def:b2:plant-life-cycles:cycles}{esporófito}.
\end{solution}

\begin{solution}{exo:b2:plant-life-cycles:5}
$0.03/10^{-4} = \qty{300}{s}$, cinco minutos. Em \qty{20}{min},
\qty{12}{cm}. A fecundação só funciona entre plantas separadas por
poucos centímetros, de modo que os \omterm{def:b2:plant-life-cycles:moss}{musgos} crescem em almofadas densas,
com os dois sexos (ou os dois órgãos) ao alcance.
\end{solution}

\begin{solution}{exo:b2:plant-life-cycles:6}
$v_s = 2\times(1.5\times 10^{-5})^2\times 1099\times 9.81/(9\times
1.8\times 10^{-5}) = \qty{0.030}{m/s}$, \qty{3}{cm/s}; distância
$uh/v_s
= 2\times 0.5/0.03 = \qty{33}{m}$. Semente: $2\times 20/0.8 =
\qty{50}{m}$ --- uma unidade mais pesada, liberada de mais alto e com
uma asa, vai mais ou menos tão longe.
\end{solution}

\begin{solution}{exo:b2:plant-life-cycles:7}
$200\times 40\times 64 = \num{512000}$ esporos por fronde; $5.1$
\omterm{def:b2:plant-life-cycles:fern}{prótalos}; $0.1$ \omterm{def:b2:plant-life-cycles:cycles}{esporófito} jovem por fronde.
\end{solution}

\begin{solution}{exo:b2:plant-life-cycles:8}
Uma semente é um \omterm{def:b2:plant-life-cycles:heterospory}{óvulo} retido: o \omterm{def:b2:plant-life-cycles:cycles}{gametófito} feminino precisa ser
mantido na planta-mãe, encerrado e alimentado, enquanto o \omterm{def:b2:plant-life-cycles:cycles}{gametófito}
masculino precisa viajar até ele. Isso exige dois tipos de esporo com
dois destinos. O único tipo de esporo de uma planta homospórica precisa
ser liberado para crescer num \omterm{def:b2:plant-life-cycles:cycles}{gametófito} livre e bissexuado; retê-lo
reteria também a função masculina e imporia a autofecundação a todas
as plantas, e não haveria nada que viajasse.
\end{solution}

\begin{solution}{exo:b2:plant-life-cycles:9}
Esporo: $\tfrac43\pi(1.5\times 10^{-5})^3\times 1100 =
\qty{1.6e-11}{kg}$, \qty{16}{ng}; energia $0.5\times 1.6\times
10^{-8}\times 2\times 10^{4} = \qty{1.6e-4}{J}$. Semente: \qty{5}{mg},
$0.5\times 5\times 10^{-3}\times 2\times 10^{4} = \qty{50}{J}$ ---
$3\times 10^{5}$ vezes mais. O esporo consegue formar uma pequena
lâmina fotossintética e nada antes de receber luz; a semente consegue
formar uma raiz e um caule no escuro, esperar durante uma estação e
sobreviver a ser comida ou ressecada.
\end{solution}

\begin{solution}{exo:b2:plant-life-cycles:10}
$10^{11}/(10^{6}\times 20) = \num{5000}$ grãos por metro cúbico; fluxo
$5000\times 10^{-7}\times 2 = \qty{1e-3}{}$ grãos por segundo; cerca de
\qty{1000}{s}, um quarto de hora, por grão. O cone precisa estar aberto
e receptivo durante horas ou dias, e exatamente quando os cones
masculinos liberam o \omterm{def:b2:plant-life-cycles:heterospory}{pólen} --- a polinização é sincronizada na escala
da semana.
\end{solution}

\begin{solution}{exo:b2:plant-life-cycles:11}
A diploidia mascara as \omterm{def:b2:genome-diversification:mutation}{mutações} recessivas deletérias; um corpo
\omterm{def:b2:meiosis-heredity:meiosis}{diploide} com tecido vascular pode ser grande e longevo, e a altura
favorece tanto a captação de luz quanto a \omterm{thm:b2:plant-life-cycles:dispersal}{dispersão dos esporos}; os
esporos produzidos no alto de um \omterm{def:b2:plant-life-cycles:cycles}{esporófito} se dispersam pelo ar, ao
passo que os gametas precisam nadar. A geração \omterm{def:b2:meiosis-heredity:meiosis}{haploide} persiste
porque a \omterm{def:b2:meiosis-heredity:meiosis}{meiose} e a fecundação exigem uma fase \omterm{def:b2:meiosis-heredity:meiosis}{haploide}, e o grão de
\omterm{def:b2:plant-life-cycles:heterospory}{pólen} e o saco embrionário são os corpos mínimos que levam os gametas
um até o outro e alimentam o embrião.
\end{solution}

\begin{solution}{exo:b2:plant-life-cycles:12}
Cultivando esporos e acompanhando-os ao microscópio, ele podia
estabelecer que um esporo cresce num pequeno corpo portador de órgãos
sexuais, que o embrião nasce da oosfera fecundada dentro do
\omterm{def:b2:plant-life-cycles:moss}{arquegônio}, que a planta portadora de esporos cresce a partir dele, e
que os mesmos órgãos existem em forma reduzida no \omterm{def:b2:plant-life-cycles:heterospory}{óvulo} das coníferas:
a alternância como uma sequência de corpos e de órgãos. Ele não podia
saber o que distinguia as duas gerações no nível da célula. As
contagens de cromossomos mostraram que os dois corpos diferem por um
fator dois no número de cromossomos, que a \omterm{def:b2:meiosis-heredity:meiosis}{meiose} ocorre na formação
dos esporos e a fecundação na oosfera, e explicaram assim por que a
alternância é obrigatória.
\end{solution}

\begin{solution}{pb:b2:plant-life-cycles:1}
\textbf{1.} $10\times 200\times 40\times 64 = 5.1\times 10^{6}$
esporos.
\textbf{2.} $\tfrac43\pi(1.5\times 10^{-5})^3\times 1100 =
\qty{1.6e-11}{kg}$ por esporo; produção $8\times 10^{-5}$ kg,
\qty{80}{mg}.
\textbf{3.} $v_s = 2\times 2.25\times 10^{-10}\times 1099\times
9.81/(1.62\times 10^{-4}) = \qty{0.030}{m/s}$.
\textbf{4.} $x = 1.5\times 0.5/0.030 = \qty{25}{m}$; disco de
$\pi\times 25^2 = \qty{2000}{m^2}$; cerca de \num{2600} esporos por
metro quadrado.
\textbf{5.} Correntes ascendentes de alguns centímetros por segundo
superam $v_s$, de modo que um esporo apanhado pela turbulência fica no
ar durante horas e percorre centenas de quilômetros; um único esporo
funda uma população se seu \omterm{def:b2:plant-life-cycles:fern}{prótalo} puder se autofecundar (ele tem os
dois órgãos), e é por isso que as ilhas remotas têm floras de
\omterm{def:b2:plant-life-cycles:fern}{samambaias} ricas.
\textbf{6.} Esporo: $0.5\times 1.6\times 10^{-8}\times 2\times 10^{4} =
\qty{1.6e-4}{J}$; produção: \qty{820}{J}; uma semente de pinheiro:
\qty{50}{J} --- a produção inteira de esporos equivale a dezesseis
sementes.
\textbf{7.} $5.1\times 10^{6}/2\times 10^{4} = 256$ \omterm{def:b2:plant-life-cycles:fern}{prótalos}.
\textbf{8.} $1 - 0.8^{6} = 1 - 0.26 = 0.74$.
\textbf{9.} $256\times 0.74\times 0.5\times 0.02 = 1.9$ \omterm{def:b2:plant-life-cycles:fern}{samambaias}
adultas por estação.
\textbf{10.} $30\times 1.9 \approx 57$: muito mais que a única
substituição de uma população estável, de modo que os valores de
sobrevivência são otimistas --- num hábitat saturado, a maioria dos
\omterm{def:b2:plant-life-cycles:cycles}{esporófitos} jovens morre por adensamento e sombreamento.
\textbf{11.} $10^{-3}/10^{-4} = \qty{10}{s}$. Muitas \omterm{def:b2:plant-life-cycles:fern}{samambaias}
amadurecem os \omterm{def:b2:plant-life-cycles:moss}{anterídios} e os \omterm{def:b2:plant-life-cycles:moss}{arquegônios} em momentos diferentes, ou
liberam um hormônio (o anteridiógeno) que torna masculinos os \omterm{def:b2:plant-life-cycles:fern}{prótalos}
jovens vizinhos, de modo que os anterozoides venham de outro
indivíduo.
\textbf{12.} Ele tem uma célula de espessura, não tem raízes nem
proteção, precisa de água líquida para a fecundação e vive poucas
semanas: quase toda a mortalidade do ciclo recai sobre ele (um esporo
em $2\times 10^{4}$ se torna um \omterm{def:b2:plant-life-cycles:fern}{prótalo}; um quarto destes nunca tem
uma noite úmida).
\textbf{13.} $v_s = 2\times 6.25\times 10^{-10}\times 599\times
9.81/(1.62\times 10^{-4}) = \qty{0.045}{m/s}$.
\textbf{14.} $3\times 15/0.045 = \qty{1000}{m}$.
\textbf{15.} $10^{11}/(2\times 10^{7}) = \num{5000}$ por metro cúbico.
\textbf{16.} $5000\times 10^{-7}\times 2 = \qty{1e-3}{s^{-1}}$: 3.6 por
hora; o primeiro chega após cerca de \qty{1000}{s}, um quarto de hora.
\textbf{17.} Cinco quilômetros a sotavento, a nuvem se espalhou
lateralmente e para cima, e a maioria dos grãos já sedimentou (um
alcance de \qty{1}{km} para cada \qty{15}{m} de altura), de modo que a
concentração é ordens de grandeza menor e um \omterm{def:b2:plant-life-cycles:heterospory}{óvulo} pode esperar dias
por um grão; o próprio \omterm{def:b2:plant-life-cycles:heterospory}{pólen} da árvore dá sobretudo sementes de
autofecundação, que abortam. A polinização pelo vento exige uma nuvem
densa; daí as populações densas.
\textbf{18.} $10^{14}$ grãos para $2\times 10^{6}$ \omterm{def:b2:plant-life-cycles:heterospory}{óvulos}: $5\times
10^{7}$ grãos por \omterm{def:b2:plant-life-cycles:heterospory}{óvulo}.
\textbf{19.} $2\times 10^{6}$ \omterm{def:b2:plant-life-cycles:heterospory}{óvulos} $\times 0.6\times 0.8 = 9.6\times
10^{5}$ sementes; \qty{4.8}{kg}; \qty{48}{MJ}.
\textbf{20.} $3\times 20/0.8 = \qty{75}{m}$: o \omterm{def:b2:plant-life-cycles:heterospory}{pólen} vai a um
quilômetro, a semente a poucas alturas de árvore; os genes viajam pelo
\omterm{def:b2:plant-life-cycles:heterospory}{pólen}, a planta pela semente.
\textbf{21.} $9.6\times 10^{5}\times 0.05\times 0.01 = 480$ descendentes
adultos, contra 57 da \omterm{def:b2:plant-life-cycles:fern}{samambaia}.
\textbf{22.} Pinheiro: $4.8\times 10^{7}/480 = \qty{100}{kJ}$ por
descendente adulto; \omterm{def:b2:plant-life-cycles:fern}{samambaia}: $820/1.9 = \qty{430}{J}$ por
descendente adulto, duzentas vezes menos.
\textbf{23.} Ambos têm metade de matéria seca a \qty{20}{kJ/g}, de
modo que, a uma taxa basal de, digamos, \qty{0.5}{mW} por grama seco,
ambos duram $10^{4}\,
\text{J/g}/(5\times 10^{-4}\,\text{W/g}) = 2\times 10^{7}$ s, cerca de
oito meses: a duração da dormência não é o que a energia da semente
compra. Ela compra o que acontece depois da germinação --- uma raiz e
um caule construídos no escuro a partir das reservas, semanas antes de
a plântula se alimentar sozinha --- e, com o envoltório, a proteção
enquanto isso; os \qty{1.6e-4}{J} do esporo constroem algumas células
que precisam fazer fotossíntese imediatamente.
\textbf{24.} Semente: fecundação sem água; um embrião abastecido,
protegido e dormente, que se estabelece em lugares secos ou sazonais;
dispersão por asas, frutos e animais. A estratégia da \omterm{def:b2:plant-life-cycles:fern}{samambaia} vence
em hábitats úmidos, sombreados e estáveis, e na colonização de
terrenos distantes ou recém-abertos, onde propágulos baratos, leves e
numerosos importam mais que o abastecimento.
\textbf{25.} \omterm{def:b2:plant-life-cycles:fern}{Samambaia}: $5.1\times 10^{6}$ esporos por estação para 1.9
adultos, $2.7\times 10^{6}$ esporos por descendente adulto,
\qty{430}{J} cada; pinheiro: $9.6\times 10^{5}$ sementes para 480
adultos, \num{2000} sementes por descendente adulto, \qty{100}{kJ}
cada.
\end{solution}
